\section{Preliminaries}
\label{sec:prelim}

We use the conventions of the preceding paper~\citep{FinCon26} and cite
its results by the numbers of its version~5, as in
\citep[Lemma~2.4]{FinCon26}. This section recalls what we use.

\subsection{Tensors and restrictions}
\label{subsec:prelim-tensors}

A \emph{tensor} over a field \(\mathbb F\) is a trilinear form in three
disjoint groups of variables, its \emph{legs}; the variables of legs
\(1\), \(2\) and \(3\) are named with the letters \(x\), \(y\) and
\(z\). We write \(T\restr T'\), and call \(T'\) a \emph{restriction} of
\(T\), if linear substitutions on each leg transform \(T\) into \(T'\):
every variable of \(T\) is replaced by a linear form in the variables of
the same leg of \(T'\). Reading a tensor as a bilinear map, with inputs on
legs \(1\) and \(2\) and the output read off by leg \(3\),
\(T\restr T'\) says that \(T'\) can be evaluated by linear maps on the
inputs, one evaluation of \(T\), and a linear map on the output. We
write \(T\cong T'\) if a bijection between the variables of each leg
transforms \(T\) into \(T'\).
A restriction or a degeneration (Section~\ref{subsec:prelim-degenerations})
\emph{holds over \(\Q\)} if all its maps have rational coefficients; it
then holds over every field of characteristic zero.

The tensor \(mT\) is the full direct sum of \(m\) copies of \(T\), with
independent variables on all three legs. Juxtaposition denotes the tensor
product, and \(T^{\otimes N}\) a tensor power. Restrictions compose, and
they are preserved by direct sums and tensor products, in particular by
tensoring with a fixed tensor. Moreover \((mT)T'\cong m(TT')\), and
reordering the factors of a tensor product is an isomorphism. For
\(\pi\in S_3\), the tensor \(T^\pi\) places leg \(i\) of \(T\) at
position \(\pi(i)\), so that \((T^{\pi'})^\pi=T^{\pi\pi'}\). In a product
of oriented factors, leg \(i\) of \(T\), in a factor \(T^\pi\), lies in
leg \(\pi(i)\) of the product, its \emph{global} leg. We write
\(\mathcal A_+\) and \(\mathcal A_-\) for the even and odd permutations
and \((23)\) for the transposition of legs \(2\) and \(3\)
\citep[Section~2.1]{FinCon26}.

\paragraph{Matrix multiplication and convolution tensors.}
The matrix multiplication tensor of side \(d\) is
\(\MM_d=\sum_{\alpha,\beta,\gamma<d}x_{\alpha\beta}y_{\beta\gamma}z_{\gamma\alpha}\),
and the dot product \(B_X(d)=x\sum_{\gamma<d}y_\gamma z_\gamma\) has a
single variable on its first leg. A sum of \(m\) products of linear forms
equal to \(T\) is the same as a restriction \(m\,\MM_1\restr T\), since
\(\MM_1\) is the single product \(xyz\). By
\citep[(2.1) and (2.2)]{FinCon26},
\begin{equation}\label{eq:prelim-mm}
 \MM_d\MM_{d'}\cong\MM_{dd'},\qquad
 \bigotimes_{\pi\in S_3}B_X(d)^\pi\cong\MM_{d^2}.
\end{equation}
For positive integers \(r\) and \(b\), the convolution tensor
\citep[(2.3)]{FinCon26} is
\[
 C(r,b)=\sum_{\alpha=0}^{r-1}\sum_{\beta=0}^{b-1}x_\alpha y_\beta z_{\alpha+\beta}.
\]
As a bilinear map it multiplies a polynomial with \(r\) coefficients
\(x_\alpha\) by a polynomial with \(b\) coefficients \(y_\beta\), and
\(z_\gamma\) reads the coefficient of the \(\gamma\)-th power in the
product. In particular \(C(1,b)=B_X(b)\). Over \(\Q\), evaluation at
the nodes \(1,\ldots,r+b-1\) followed by Lagrange interpolation writes
\(C(r,b)\) as a sum of \(r+b-1\) products of linear forms with rational
coefficients \citep[Lemma~9.1(a)]{FinCon26}. The symmetrized tensors
\citep[(2.5) and (2.7)]{FinCon26} are
\[
 S^\sigma_{r,b}=\bigotimes_{\pi\in\mathcal A_\sigma}C(r,b)^\pi
 \quad(\sigma\in\{+,-\}),\qquad
 Q_{r,b}=S^+_{r,b}S^-_{r,b}=\bigotimes_{\pi\in S_3}C(r,b)^\pi .
\]
By \citep[(2.6)]{FinCon26}, \(S^\sigma_{r,b}\cong S^{-\sigma}_{b,r}\)
and \(S^\sigma_{1,b}\cong\MM_b\). With~\eqref{eq:prelim-mm} and six
oriented copies of the rational decomposition, this gives (the first
identity is also \citep[(2.8)]{FinCon26})
\begin{equation}\label{eq:prelim-Q}
 Q_{r,b}\cong Q_{b,r},\qquad Q_{1,t}\cong\MM_{t^2},\qquad
 (r+b-1)^6\,\MM_1\restr Q_{r,b}\ \text{over }\Q .
\end{equation}

\subsection{Degenerations and exactification}
\label{subsec:prelim-degenerations}

Let \(\theta\ge0\) be an integer. We write \(T\degen T'\), a
\emph{degeneration of degree at most \(\theta\)}, if a triple \(\mathsf S(t)\) of
linear maps, one on each leg, with Laurent polynomials in \(t\) as
entries, transforms \(T\) into a polynomial in \(t\) of degree at most
\(\theta\) with tensor coefficients and constant term \(T'\)
\citep[Section~2.3]{FinCon26}. The entries of \(\mathsf S(t)\) may contain
negative powers of \(t\); the transformed tensor may not. Outside
degenerations, \(t\) denotes the positive integer in the terminal pair
\((1,t)\) of a chain (Section~\ref{subsec:chain-counts}).

We use two facts. First, \citep[Lemma~2.1]{FinCon26} (Composition): maps
that do not depend on \(t\), applied before a degeneration, or after it if
they are defined on all its target variables, keep the degree bound and
act on the constant term; applying degenerations to the factors of a
tensor product with one common parameter \(t\) gives a degeneration of
the product whose degree bound is the sum of the degree bounds and whose
constant term is the product of the constant terms; and both statements
hold for oriented copies \(T^\pi\), with the maps moved to the permuted
legs. Second, \citep[Lemma~2.2]{FinCon26} (Exactification, after
\citet{Bini80}): if \(T\degen T'\) with degree at most \(\theta\), then
\[
 (k\theta+1)\,T^{\otimes k}\restr T'^{\otimes k}\qquad(k\ge1).
\]
That lemma uses roots of unity. By \citep[Lemma~9.1(c)]{FinCon26}, the same
restriction is given by the maps \(\mathsf S(\ell)^{\otimes k}\),
\(1\le\ell\le k\theta+1\), with the scalars \((-1)^{\ell-1}\binom{k\theta+1}\ell\)
applied on one fixed leg; so it holds over \(\Q\) when \(\mathsf S\) has
rational coefficients.

Lemma~2.1 of \citep{FinCon26} does not cover two degenerations applied
in sequence: the Laurent maps of the second would act on the terms of
positive degree of the first and could create negative powers of \(t\).
So the chains of Section~\ref{subsec:chain-counts} exactify each stage
before its output is used.

\subsection{Relative schemes}
\label{subsec:prelim-schemes}

A \emph{relative scheme} \(r_*\MM_A\restr m_*\MM_B\), with positive
integers \(r_*\), \(m_*\) and \(A<B\), is as in
Section~\ref{subsec:intro-results} \citep[Definition~2.3]{FinCon26}: it
is given by its three linear maps, and has source side \(A\), target
side \(B\) and pooled exponent \(p=\log(r_*/m_*)/\log(B/A)\). By the
pooled recursion \citep[Lemma~2.4]{FinCon26}, \(p\ge2\), and the scheme
yields a fixed arithmetic algorithm that multiplies \(n\times n\)
matrices with \(O(n^p)\) operations if \(p>2\), and with
\(O(n^2\log n)\) operations if \(p=2\). The algorithm uses only the
coefficients of the scheme, so in both cases \(\omega\le p\) over every
field that contains them. We also use the \emph{fixed-batch bound}
of \citep[(10.1)]{FinCon26},
\begin{equation}\label{eq:prelim-phi}
 \Phi=\frac{\log(2r_*/m_*)}{\log(B/A)}=p+\frac{\log2}{\log(B/A)},
\end{equation}
which exceeds \(p\). As \(p>0\) gives \(r_*/m_*>1\), and hence
\(\lceil r_*/m_*\rceil<2r_*/m_*\), it also bounds the exponent
\(\log\lceil r_*/m_*\rceil/\log(B/A)\) of the recursion without pooling
\citep[Section~10.1]{FinCon26}.

\subsection{Sectors, type restrictions and the separator}
\label{subsec:prelim-sectors}

Following \citep[Section~4]{FinCon26}, a tensor
\(T=\sum_{g=0}^{M-1}T_g\) with \(T_g\in X\otimes Y_g\otimes Z_g\),
\(Y=\bigoplus_gY_g\) and \(Z=\bigoplus_gZ_g\) is a sum of \(M\)
\emph{sectors sharing the first leg}. It is not a direct sum: all sectors
use the same first-leg variables, so a coordinate change that depends on
the sector can be applied on legs \(2\) and \(3\) but not yet on leg~\(1\).

Sectors arise from \emph{type restrictions} \citep[Section~4.1 and
Lemma~4.3]{FinCon26}, which keep in a tensor power only the coordinates
on legs \(2\) and \(3\) whose words of labels have a prescribed type;
Lemma~\ref{lem:stage-types} adapts the proof of
\citep[Lemma~4.3]{FinCon26}, with weights on classes of labels in place
of the filtration index.

The \emph{separator} gives each sector its own first leg. Shared-leg
separation is due to \citet[Proposition~3.1]{OAI26}; we use the
polynomial separator of \citet[Proposition~2.1]{Lu26}, in the rational
form \citep[Lemma~9.1(b)]{FinCon26}, which writes it in the coordinates
of \citep[Proposition~4.6]{FinCon26} (see \citep[Remark~4.7]{FinCon26}).
For \(T\) as above, it is a degeneration
\begin{equation}\label{eq:prelim-separator}
 (2M-1)\,T\degen\bigoplus_{g=0}^{M-1}T_g\,B_X(M)\qquad\text{over }\Q,
\end{equation}
of degree at most \(M-1\), in which each summand on the right has its own
copy of \(X\); permuting legs gives the same with any shared leg. The
target variables are \(X_{\alpha,g'}\), \(Y_{g,\beta,e}\) and
\(Z_{g,\gamma,e'}\) with \(g',e,e'<M\), where \(\alpha\), \(\beta\)
and \(\gamma\) index coordinates of \(X\), \(Y_g\) and \(Z_g\); the
constant term consists of the terms with \(g'=g\) and \(e=e'\). Suppose
that coordinate permutations \((\rho^X_g,\rho^Y_g,\rho^Z_g)\) identify
each sector \(T_g\) with one fixed tensor \(\mathcal V\). Apply \(\rho^X_{g'}\) to
the copy \((X_{\alpha,g'})_\alpha\), and \(\rho^Y_g\) and \(\rho^Z_g\) to
the variables \(Y_{g,\beta,e}\) and \(Z_{g,\gamma,e'}\), for all
\(g,g',e,e'\). This \emph{reindexing} \citep[(B.6)]{FinCon26} is a
bijection of all target variables that does not depend on \(t\), so by
\citep[Lemma~2.1]{FinCon26} it can follow the separator; it maps the
constant term of~\eqref{eq:prelim-separator} onto
\(M\,(\mathcal V B_X(M))\).

\subsection{Profiles}
\label{subsec:prelim-profiles}

A \emph{profile} \citep[Definition~3.1]{FinCon26} is a positive symmetric
function \(P\) on pairs of positive integers such that, for all positive
integers \(r\) and \(b\),
\[
 P(1,b)=b,\qquad
 2P(r,b)\ge P(r,b-1)+P(r,b+1)\ \ (b\ge2),\qquad
 P(r,3b+r-1)\ge3P(r,b).
\]
These are the normalization, concavity and tripling axioms
\citep[(3.1)--(3.3)]{FinCon26}. The \emph{least profile} \(F\)
of~\citep[(3.4)]{FinCon26} is a profile, and every profile \(P\) satisfies
\(P\ge F\) at every pair \citep[Theorem~3.2]{FinCon26}.
Section~\ref{subsec:chains-ramp} recalls the formula for \(F\), and
Lemma~\ref{lem:profile-potential} uses \(F\) as a potential.
